% =============================================================================
% proofs/05_existence.tex
% Cutoff equilibrium at order size two, conditional on a certified box.
% Statements first, then the proofs.  Assembled into appendix.tex.
% Needs amsmath, amssymb, amsthm and enumitem in the assembling preamble.
% =============================================================================

\providecommand{\Eop}{\mathbb{E}}
\providecommand{\Prb}{\mathbb{P}}
\providecommand{\ind}{\mathbf{1}}
\providecommand{\Tmap}{\mathsf{T}}

\makeatletter
\@ifundefined{c@theorem}{\newtheorem{theorem}{Theorem}}{}
\@ifundefined{c@lemma}{\newtheorem{lemma}{Lemma}}{}
\@ifundefined{c@remark}{\newtheorem{remark}{Remark}}{}
\makeatother

\section{Cutoff equilibrium at order size two}
\label{sec:pf-existence}

\subsection*{Setting and notation}
\label{sec:pf-existence-setup}

Trading occupies dates $d=0,\ldots,H$. After date $H$ a bidder observes the public information set
and enters with probability $p(\mathcal I_H)$. That is the two-round game: a trading round, then a
control round. Noise takes the values $-\bar z$, $0$ and $+\bar z$ with probabilities $\kappa/2$,
$1-\kappa$ and $\kappa/2$. On a building date the engaged blockholder's order is two noise lumps, so
pooled order flow takes five values. The plan menu is Exit, Hold and one Voice family. Exit sells
the initial stake $b_0$ at date $0$ and is idle thereafter; Hold is idle at $b_0$ throughout; Voice
accumulates from $b_0$ to a strictly increasing terminal target $b^*(s)$ in $n(s)$ building dates,
with $n$ a weakly decreasing step function of the signal. Only Voice ever crosses the threshold
$\tau$, and the legal clock reads the crossing off the stake path. There is no feedback from
realised order flow into the executed path.

A cutoff policy is a pair $k=(k_1,k_2)$ in the ordered polytope
\begin{equation}
\label{eq:ex-theta}
\Theta \;:=\; \bigl\{\, (k_1,k_2):\ s_{\mathrm{lo}}\le k_1\le k_2\le s_{\mathrm{hi}} \,\bigr\},
\end{equation}
with Exit on $[s_{\mathrm{lo}},k_1)$, Hold on $[k_1,k_2)$ and Voice on $[k_2,s_{\mathrm{hi}}]$. The
signal is taken on the compact interval $[s_{\mathrm{lo}},s_{\mathrm{hi}}]$ used at the calibration,
with the Gaussian density truncated and renormalised. Write $U_j(s;k)$ for the expected payoff of
plan $j$ at signal $s$ under the price system that $k$ induces, and write the adjacent-plan gaps
\begin{equation}
\label{eq:ex-gaps}
g^{\mathrm{EH}}(s;k) \;:=\; U_{\mathrm{Exit}}(s;k)-U_{\mathrm{Hold}}(s;k),
\qquad
g^{\mathrm{HV}}(s;k) \;:=\; U_{\mathrm{Hold}}(s;k)-U_{\mathrm{Voice}}(s;k).
\end{equation}
Write $\mathcal S_{\mathrm{free}}$ for the finite set of $k$-free breakpoints in
$(s_{\mathrm{lo}},s_{\mathrm{hi}})$: the jump points of $n(s)$, together with the signals at which
the Voice stake path first reaches $\tau$ on some date. Equivalently, $\mathcal S_{\mathrm{free}}$
is the breakpoint set of the menu computed with both cutoffs placed outside the support. These
points do not move with $k$. A \emph{down-crossing} of a gap $g(\cdot;k)$ is a point
$z\in[s_{\mathrm{lo}},s_{\mathrm{hi}}]$ at which $g(s_{\mathrm{lo}};k)\ge 0\ge g(s_{\mathrm{hi}};k)$,
$g(s;k)\ge 0$ for every $s<z$ off $\mathcal S_{\mathrm{free}}$, and $g(s;k)\le 0$ for every $s>z$
off $\mathcal S_{\mathrm{free}}$. The down-crossing is unique when no other point of the interval
satisfies those three properties. The $k$-free set is finite, so this is uniqueness of the sign
change of a function that is continuous in $s$ on the complementary open intervals.

The inner price at a public history solves $P=\Eop[Y\mid\mathcal I]$, equivalently
\begin{equation}
\label{eq:ex-inner}
P \;=\; \bigl(1-p(P)\bigr)\hat V \;+\; p(P)\bigl(P+\tilde m\bigr),
\qquad
p(P) \;=\; 1-\Phi\!\left(\frac{P+K+\tilde m-\bar S}{\sigma_\xi}\right),
\end{equation}
with $\hat V=\hat v+\pi\Delta_V$ and $\tilde m=m_0+\pi\Delta_m$ the belief summaries at that
history. Off-path pooled histories, those that no type selected by $k$ can produce, are priced from
a $k$-independent reference: each empty mark path is assigned the mean fundamental of the $n(s)$
cell that would produce it. Flagged histories are priced from the filing $(B^F,Q^F,a{=}1)$ alone.
The Voice target is
\begin{equation}
\label{eq:ex-bstar}
b^*(s) \;=\; b_0+(\bar b-b_0)\cdot\tfrac12\Bigl(1+\frac{x}{\sqrt{1+x^2}}\Bigr),
\qquad x=(s-\mu_v)/\sigma_s,
\end{equation}
so $b^{*\prime}(s)=\tfrac12(\bar b-b_0)(1+x^2)^{-3/2}/\sigma_s>0$ on the whole line whenever
$\bar b>b_0$ and $\sigma_s>0$. On the flagged cell $B^F+Q^F=b^*(s)$, and the strictly increasing
target therefore pins the signal from the filing. The flagged price at a signal is consequently
well-defined independently of $k$.

Given $k\in\Theta$ at which each gap in \eqref{eq:ex-gaps} has a unique down-crossing, write
$z_{\mathrm{EH}}(k)$ and $z_{\mathrm{HV}}(k)$ for those two points and set
\begin{equation}
\label{eq:ex-T}
\Tmap(k) \;:=\; \bigl(z_{\mathrm{EH}}(k),\, z_{\mathrm{HV}}(k)\bigr)
\end{equation}
when the zeros are ordered $z_{\mathrm{EH}}(k)\le z_{\mathrm{HV}}(k)$. That is the outer map. An
\emph{equilibrium} of the two-round game is a cutoff $k^\star\in\Theta$ at which, under the unique
inner price system induced by $k^\star$ and the on-path Bayes and off-path reference beliefs just
named, the plan that $k^\star$ selects maximises $U_j(s;k^\star)$ at every signal $s$. Equivalently,
$k^\star=\Tmap(k^\star)$, once the ordered unique down-crossings make adjacent indifference global
(Lemma~\ref{lem:ex-global} below).

\subsection*{Conditions}
\label{sec:pf-existence-conditions}

Let $B=[k_1^-,k_1^+]\times[k_2^-,k_2^+]$ be a compact axis-aligned box in the cutoff plane. The four
conditions below are named, and the theorem is conditional on all four.

\begin{enumerate}[label=(B\arabic*),leftmargin=2.6em,itemsep=4pt]
\item \label{cond:B1} \emph{Interior ordered box.} The box sits in the interior of $\Theta$ and is
  strictly ordered:
  \[
  s_{\mathrm{lo}} \;<\; k_1^- \;\le\; k_1^+ \;<\; k_2^- \;\le\; k_2^+ \;<\; s_{\mathrm{hi}}.
  \]
\item \label{cond:B2} \emph{Breakpoint-free Voice coordinate.} The Voice interval of the box is
  disjoint from the $k$-free set: $[k_2^-,k_2^+]\cap\mathcal S_{\mathrm{free}}=\emptyset$.
\item \label{cond:B3} \emph{Unique ordered down-crossings.} At every $k\in B$, each of
  $g^{\mathrm{EH}}(\cdot;k)$ and $g^{\mathrm{HV}}(\cdot;k)$ has a unique down-crossing on
  $[s_{\mathrm{lo}},s_{\mathrm{hi}}]$, and the zeros are ordered
  $z_{\mathrm{EH}}(k)\le z_{\mathrm{HV}}(k)$.
\item \label{cond:B4} \emph{Miranda face signs.} Writing $\Tmap=(\Tmap_1,\Tmap_2)$, the displacement
  $F(k):=\Tmap(k)-k$ points inward on each face of $B$:
  \begin{align}
  \label{eq:ex-miranda}
  \Tmap_1(k_1^-,k_2) &\ge k_1^-
    &&\text{for all }k_2\in[k_2^-,k_2^+],
  \nonumber\\
  \Tmap_1(k_1^+,k_2) &\le k_1^+
    &&\text{for all }k_2\in[k_2^-,k_2^+],
  \\
  \Tmap_2(k_1,k_2^-) &\ge k_2^-
    &&\text{for all }k_1\in[k_1^-,k_1^+],
  \nonumber\\
  \Tmap_2(k_1,k_2^+) &\le k_2^+
    &&\text{for all }k_1\in[k_1^-,k_1^+].
  \nonumber
  \end{align}
\end{enumerate}

\subsection*{Statements}
\label{sec:pf-existence-statements}

\begin{lemma}[Unique inner price]
\label{lem:ex-inner}
Assume $m_0>0$, $\Delta_m>0$ and $\sigma_\xi>0$. Then for every belief pair
$(\hat v,\pi)\in\mathbb R\times[0,1]$ the inner equation \eqref{eq:ex-inner} has a unique real root
$P$, and that root is continuous in $(\hat v,\pi)$.
\end{lemma}

\begin{lemma}[The Exit--Hold gap decreases in the signal]
\label{lem:ex-eh}
Assume $b_0>0$, $\sigma_v>0$, $\sigma_\varepsilon>0$ and the hypotheses of
Lemma~\ref{lem:ex-inner}. At every $k\in\Theta$, the gap $g^{\mathrm{EH}}(\cdot;k)$ is affine and
strictly decreasing on $[s_{\mathrm{lo}},s_{\mathrm{hi}}]$. A unique down-crossing of this gap is
therefore equivalent to the endpoint sign pattern
$g^{\mathrm{EH}}(s_{\mathrm{lo}};k)\ge 0\ge g^{\mathrm{EH}}(s_{\mathrm{hi}};k)$.
\end{lemma}

\begin{lemma}[Ordered adjacent indifference is global]
\label{lem:ex-global}
Fix $k\in\Theta$ at which both gaps have unique down-crossings $z_{\mathrm{EH}}\le z_{\mathrm{HV}}$.
Then, off the finite set $\mathcal S_{\mathrm{free}}$, Exit maximises $U_j(s;k)$ on
$[s_{\mathrm{lo}},z_{\mathrm{EH}}]$, Hold maximises it on $[z_{\mathrm{EH}},z_{\mathrm{HV}}]$, and
Voice maximises it on $[z_{\mathrm{HV}},s_{\mathrm{hi}}]$. The $k$-free set is null for the signal
law. In particular, if $k=\Tmap(k)$, the cutoff policy $k$ is an equilibrium.
\end{lemma}

\begin{theorem}[Cutoff equilibrium on a certified box]
\label{thm:existence}
Fix a disclosure rule $(\tau,T)$ and a liquidity $\kappa\in(0,1)$. Assume $m_0>0$, $\Delta_m>0$,
$\sigma_\xi>0$, $\sigma_v>0$, $\sigma_\varepsilon>0$, and $0<b_0<\tau<\bar b$. Let
$B=[k_1^-,k_1^+]\times[k_2^-,k_2^+]$ be a compact box, and assume (B1)--(B4). Then the outer map
$\Tmap$ is well-defined and continuous on $B$, and there exists $k^\star\in B$ with
$\Tmap(k^\star)=k^\star$. The cutoff $k^\star$, together with the unique inner price system it
induces, is an equilibrium of the two-round game at order size two.
\end{theorem}

\begin{remark}[What a probe grid does and does not show]
\label{rem:ex-probe}
The numerical record in \texttt{numerical\_v4/checks/t5\_existence\_conditions.json} builds, at each
calibration node, a box $B$ about the solver's candidate. It checks (B1) and (B2) on that box, and
it evaluates (B3) and (B4) on the stated $3\times 3$ probe grid on $B$ (the four vertices, the four
face midpoints, and the centre), together with the inner-root certificate of
Lemma~\ref{lem:ex-inner} and the flagged-target slope at those nine points. A calibration node is a
pair $(T,\tau_q)$ with $T\in\{5,10\}$ and $\tau_q$ one of the five frozen quantiles
$(0.1,0.3,0.5,0.7,0.9)$ of the seed equilibrium's Voice $b^*(s)$ distribution, at baseline
$\kappa=0.5$. At each probe the gaps are evaluated on a $2001$-point signal grid. The probe grid is
not a proof of (B3) at every point of $B$, nor of (B4) at every point of each face.
Theorem~\ref{thm:existence} is conditional on (B1)--(B4).
\end{remark}

\subsection*{Proofs}
\label{sec:pf-existence-proofs}

\begin{proof}[Proof of Lemma~\ref{lem:ex-inner}]
Write $\tilde m=m_0+\pi\Delta_m$. The standing signs give $\tilde m\ge m_0>0$ for every
$\pi\in[0,1]$. The entry map $p$ is $C^\infty$ and strictly decreasing, with values in $(0,1)$ at
every finite $P$, because $\sigma_\xi>0$ and $\Phi$ does not attain $0$ or $1$ at a finite argument.
Rearrange \eqref{eq:ex-inner} as $P=\psi(P)$ with
\[
\psi(P) \;:=\; \hat V \;+\; \tilde m\,\frac{p(P)}{1-p(P)}.
\]
The ratio $p/(1-p)$ is strictly increasing in $p\in(0,1)$, hence strictly decreasing in $P$, and
$\tilde m>0$, so $\psi$ is strictly decreasing. The identity map is strictly increasing, so they
cross at most once. For existence: as $P\to-\infty$, $p(P)\to 1$ and $\psi(P)\to+\infty$, so
$\psi(P)>P$; as $P\to+\infty$, $p(P)\to 0$ and $\psi(P)\to\hat V$, so $\psi(P)<P$. Continuity of
$\psi$ on $\mathbb R$ and the two endpoint inequalities give a root, which is unique by the opposite
monotonicities. The unique root of a continuous function that is strictly increasing in $P$ (namely
$P-\psi(P)$) and jointly continuous in $(P,\hat v,\pi)$ varies continuously with $(\hat v,\pi)$.
\end{proof}

\begin{proof}[Proof of Lemma~\ref{lem:ex-eh}]
Exit and Hold share the idle mark path, so both are priced off the same pooled histories. Exit has
terminal stake $0$ and sells $b_0$ at date $0$, hence
$U_{\mathrm{Exit}}(s;k)=b_0\,\Eop[P_0^P\mid \text{idle}]$, a number that does not depend on $s$.
Hold has terminal stake $b_0$, makes no trade, and carries $a=0$, so
\[
U_{\mathrm{Hold}}(s;k)
 \;=\; b_0\Bigl(
   \bigl(1-\Eop[p_H]\bigr)\bigl(\mu_v+\beta(s-\mu_v)\bigr)
   +\Eop[p_H P_H]
   +\Eop[p_H]\,m_0
 \Bigr),
\]
with $\beta=\sigma_v^2/(\sigma_v^2+\sigma_\varepsilon^2)\in(0,1)$ and with $(p_H,P_H)$ the
control-node entry probability and price along the idle path, both independent of $s$. Lemma~\ref{lem:ex-inner}
puts every inner price at a finite value, so $p_H\in(0,1)$ and $\Eop[p_H]\in(0,1)$. Therefore
\[
g^{\mathrm{EH}}(s;k)
 \;=\; C(k) \;-\; b_0\bigl(1-\Eop[p_H]\bigr)\beta\, s
\]
for a constant $C(k)$ in $s$. The coefficient of $s$ is strictly negative, so the gap is affine and
strictly decreasing. A strictly decreasing continuous function of one variable has a unique
down-crossing on a compact interval if and only if it is nonnegative at the left endpoint and
nonpositive at the right endpoint.
\end{proof}

\begin{proof}[Proof of Lemma~\ref{lem:ex-global}]
Write $z_1=z_{\mathrm{EH}}$ and $z_2=z_{\mathrm{HV}}$, with $z_1\le z_2$ by hypothesis, and suppress
$k$. Unique down-crossing of $g^{\mathrm{EH}}$ gives $U_{\mathrm{Exit}}\ge U_{\mathrm{Hold}}$ on
$[s_{\mathrm{lo}},z_1]$ and $U_{\mathrm{Hold}}\ge U_{\mathrm{Exit}}$ on $[z_1,s_{\mathrm{hi}}]$.
Unique down-crossing of $g^{\mathrm{HV}}$ gives $U_{\mathrm{Hold}}\ge U_{\mathrm{Voice}}$ on
$[s_{\mathrm{lo}},z_2]$ and $U_{\mathrm{Voice}}\ge U_{\mathrm{Hold}}$ on $[z_2,s_{\mathrm{hi}}]$.

On $[s_{\mathrm{lo}},z_1]$ one has $s\le z_1\le z_2$, so $U_{\mathrm{Exit}}\ge U_{\mathrm{Hold}}\ge
U_{\mathrm{Voice}}$. On $[z_1,z_2]$ one has $s\ge z_1$ and $s\le z_2$, so
$U_{\mathrm{Hold}}\ge U_{\mathrm{Exit}}$ and $U_{\mathrm{Hold}}\ge U_{\mathrm{Voice}}$. On
$[z_2,s_{\mathrm{hi}}]$ one has $s\ge z_2\ge z_1$, so $U_{\mathrm{Voice}}\ge U_{\mathrm{Hold}}\ge
U_{\mathrm{Exit}}$. The three comparisons are on the complement of $\mathcal S_{\mathrm{free}}$,
where the gaps are continuous in $s$. At each such signal a maximiser is the plan the ordered
cutoffs $(z_1,z_2)$ select, with indifference between adjacent plans at the two cutoffs themselves.
If $k=\Tmap(k)$, then $(k_1,k_2)=(z_1,z_2)$ and the cutoff policy is an equilibrium.
\end{proof}

\begin{proof}[Proof of Theorem~\ref{thm:existence}]
\emph{Step 1 (inner prices).} Lemma~\ref{lem:ex-inner} supplies a unique inner price at every public
history, flagged or pooled, continuously in the belief summaries. The standing signs $m_0>0$ and
$\Delta_m>0$ make $\tilde m>0$ at every posterior engagement share, and $\sigma_\xi>0$ makes the
entry map strictly monotone, so uniqueness is not a property of $B$: it holds at every $k\in\Theta$.

\emph{Step 2 ($\Tmap$ is well-defined on $B$).} Condition (B3) gives, at every $k\in B$, unique
down-crossings of both gaps with $z_{\mathrm{EH}}(k)\le z_{\mathrm{HV}}(k)$. The right-hand side of
\eqref{eq:ex-T} is therefore a single point of $\Theta$, and $\Tmap:B\to\Theta$ is well-defined.
Lemma~\ref{lem:ex-eh} records that the Exit--Hold half of (B3) is the endpoint sign pattern of an
affine gap; the Hold--Voice half is the content of (B3) that is not supplied by the menu.

\emph{Step 3 (type masses are continuous on $B$, and the populated type set is constant).} The
$k$-free set $\mathcal S_{\mathrm{free}}$ is finite, hence closed. Condition (B2) says that the
compact interval $[k_2^-,k_2^+]$ is disjoint from it, so the two are a positive distance apart.
Condition (B1) puts $k_1^+<k_2^-$, so as $k$ ranges over $B$ the Voice region $[k_2,s_{\mathrm{hi}}]$
gains or loses only signals in $[k_2^-,k_2^+]$, none of which is an $n$-jump or a $\tau$-crossing
pull-back. The Voice types that carry positive mass are therefore the same at every $k\in B$: the
lowest Voice atom is $[k_2,s_{\mathrm{next}})$, where $s_{\mathrm{next}}$ is the next point of
$\mathcal S_{\mathrm{free}}\cup\{s_{\mathrm{hi}}\}$ above $k_2^+$, and every $n$-plateau lying
entirely above $s_{\mathrm{next}}$ is included in full. Type $0$, the idle path shared by Exit and
Hold, occupies $[s_{\mathrm{lo}},k_2)$ and has mass at least the Gaussian mass of
$[s_{\mathrm{lo}},k_2^-)$, which is positive by (B1). Types whose $n$-plateau lies entirely below
$k_2^-$ are empty throughout $B$, and their beliefs are the constant reference. Each populated
type's mass and truncated mean fundamental are Gaussian cdf and truncated-mean evaluations at
endpoints that are either fixed ($k$-free breakpoints, $s_{\mathrm{lo}}$, $s_{\mathrm{hi}}$) or
equal to $k_1$ or $k_2$. Those evaluations are continuous in $k$ on $B$, and the masses of
populated types are bounded away from zero on the compact box.

\emph{Step 4 (prices and payoffs are continuous in $k$).} At a pooled history the posterior
summaries $(\hat v,\pi)$ are ratios of linear forms in the type masses, with likelihoods that do not
depend on $k$. On histories feasible for a populated type the denominator is bounded away from zero
on $B$ by Step~3. On histories feasible only for types empty throughout $B$ the summaries equal the
$k$-independent reference. In both cases $(\hat v,\pi)$ is continuous in $k$ on $B$. Step~1 upgrades
that to continuity of every inner price. Expected execution prices along each mark path are finite
sums of those inner prices against $k$-free likelihoods, hence continuous in $k$. For each plan $j$
and each signal $s\notin\mathcal S_{\mathrm{free}}$, the stake path, the mark path, the clock and the
engagement cost are locally constant in $s$ and independent of $k$, so $U_j(s;k)$ is continuous in
$k$ on $B$. (The flagged price at a given $s$ is itself $k$-free, by injectivity of $b^*$,
and the remaining $k$-dependence is a continuous pre-filing trade cost.) Consequently both gaps in
\eqref{eq:ex-gaps} are continuous in $k$ at every $s\notin\mathcal S_{\mathrm{free}}$.

\emph{Step 5 ($\Tmap$ is continuous on $B$).} Let $k^n\to k$ in $B$ and write
$z^n=z_{\mathrm{EH}}(k^n)$. The sequence $z^n$ lives in the compact interval
$[s_{\mathrm{lo}},s_{\mathrm{hi}}]$, so a subsequence $z^{n_j}$ converges to some
$z_\infty$ in that interval. For any $s<z_\infty$ with $s\notin\mathcal S_{\mathrm{free}}$, one has
$s<z^{n_j}$ for large $j$, hence $g^{\mathrm{EH}}(s;k^{n_j})\ge 0$, and Step~4 passes to the limit
$g^{\mathrm{EH}}(s;k)\ge 0$. For any $s>z_\infty$ off $\mathcal S_{\mathrm{free}}$, the same
argument gives $g^{\mathrm{EH}}(s;k)\le 0$. The finite set $\mathcal S_{\mathrm{free}}$ is excluded from the down-crossing definition, so
$z_\infty$ is a down-crossing of $g^{\mathrm{EH}}(\cdot;k)$. Condition (B3) says that
down-crossing is unique, hence $z_\infty=z_{\mathrm{EH}}(k)$. Every convergent subsequence of $z^n$ has this same limit, so
$z_{\mathrm{EH}}(k^n)\to z_{\mathrm{EH}}(k)$. The identical argument with $g^{\mathrm{HV}}$ gives
continuity of $z_{\mathrm{HV}}$. Thus $\Tmap$ is continuous on $B$.

\emph{Step 6 (Miranda's theorem).} The box $B$ is a product of two compact intervals of positive
length, by (B1). Define $G(k):=k-\Tmap(k)$, continuous on $B$ by Step~5. Condition (B4) is exactly
the classical face-sign pattern for $G$: $G_1\le 0$ on the face $k_1=k_1^-$, $G_1\ge 0$ on the face
$k_1=k_1^+$, $G_2\le 0$ on the face $k_2=k_2^-$, and $G_2\ge 0$ on the face $k_2=k_2^+$. Miranda's
theorem (the intermediate-value theorem on a product of compact intervals) therefore supplies a
point $k^\star\in B$ with $G(k^\star)=0$, that is $\Tmap(k^\star)=k^\star$. The argument uses that
$\Tmap$ is a well-defined continuous map on the compact rectangle $B$. It does not require
$\Tmap(B)\subseteq B$.

\emph{Step 7 (the fixed point is an equilibrium).} The point $k^\star$ lies in $B\subset\Theta$. At
$k^\star$ the two unique down-crossings are the coordinates of $k^\star$ itself, and they are
ordered by (B3). Lemma~\ref{lem:ex-global} upgrades adjacent indifference to global optimality of
the selected plan at every signal off the null $k$-free set. The inner price system is unique by Step~1, on-path beliefs are
Bayes from the type masses of $k^\star$, and off-path beliefs are the model's $k$-independent
reference. This is an equilibrium of the two-round game at order size two.
\end{proof}
